\section{Compilation Request Flow}\label{sec:request-flow}

The compile loop is the central integration contract between the Monaco editor, the Express API, and the native
compiler.
It is intentionally narrow: the client sends UTF-8 source text; the server returns either MIDI bytes or a JSON array
of diagnostics; the compiler remains a CLI subprocess with a stable interface.

\subsection{End-to-end sequence}\label{subsec:compile-sequence}

\begin{figure}[htbp]
  \centering
  \motivoscalefig{%
    \begin{tikzpicture}[
        actor/.style={motivo/actor, minimum width=2.2cm, minimum height=0.65cm},
        msg/.style={motivo/msg},
        arrow/.style={motivo/arrow}
      ]
      \node[actor, fill=motivoBlue!12] (editor) at (0,0) {Editor};
      \node[actor, fill=motivoBlue!10] (client) at (2.9,0) {Compile client};
      \node[actor, fill=motivoCyan!15] (api) at (5.8,0) {Express};
      \node[actor, fill=motivoGreen!15] (bin) at (8.7,0) {\texttt{motivoc}};
      \node[actor, fill=motivoAmber!18] (ui) at (11.6,0) {Visualizer};

      \foreach \x in {0,2.9,5.8,8.7,11.6} {
        \draw[draw=motivoInk!30] (\x,-0.45) -- (\x,-4.9);
      }

      \draw[arrow] (0,-0.85) -- node[msg, above] {source text} (2.9,-0.85);
      \draw[arrow] (2.9,-1.45) -- node[msg, above] {\texttt{POST /api/compile}} (5.8,-1.45);
      \draw[arrow] (5.8,-2.05) -- node[msg, above] {temp \texttt{.motivo}} (8.7,-2.05);
      \draw[arrow] (8.7,-2.65) -- node[msg, above] {stderr or \texttt{.mid}} (5.8,-2.65);
      \draw[arrow] (5.8,-3.25) -- node[msg, above] {MIDI or JSON} (2.9,-3.25);
      \draw[arrow] (2.9,-3.85) -- node[msg, above] {bytes / diagnostics} (0,-3.85);
      \draw[arrow] (0,-4.45) -- node[msg, above] {markers + notes} (11.6,-4.45);
    \end{tikzpicture}%
  }
  \caption{Compile request sequence from editor action through \texttt{/api/compile} to visualization.
  Next.js rewrites \texttt{/api/compile} to the Express \texttt{POST /compile} handler.}\label{fig:compile-sequence}
\end{figure}

When the user triggers compile (toolbar control or shortcut), the client reads the current editor buffer and posts it
to the compile endpoint.
The HTTP client posts \texttt{\{ source \}} to \texttt{/api/compile} with \texttt{responseType: arraybuffer} so a
successful compile can return raw MIDI without an extra download step.

On the server, Zod validates the body, then the compiler service writes source and output to temporary files, runs
\texttt{motivoc}, and deletes both files when the request completes.
Non-zero exit codes with stderr content are parsed line-by-line into diagnostics (lexical, syntax, semantic, lowering,
output stages) using a regex that strips ANSI color codes and optional \texttt{line:column} locations; each parsed
record exposes a \texttt{type} field in the JSON API (the compiler stderr labels name the pipeline stage).
Infrastructure failures (missing binary, I/O errors, missing output file) map to HTTP~500 responses with stable error
codes for the client.

On success, the controller sets \texttt{Content-Type: audio/midi} and sends the file bytes.
On semantic or syntax failure, the controller responds with HTTP~400 and a JSON body \texttt{\{ kind: 'error',
diagnostics: [...] \}}.

The client clears previous MIDI and editor markers, then either loads bytes into \texttt{MidiContext} for piano-roll
rendering and playback, or converts diagnostics to Monaco markers via \texttt{setError} and \texttt{jumpTo} on the first
error, while mirroring the full list in the logs panel.

\subsection{Testability of the contract}\label{subsec:compile-contract-tests}

This boundary is tested at three levels without always building C++.
Server Vitest suites inject fake compiler executables to assert success paths, stderr parsing, and infrastructure
errors.
Client tests cover diagnostic mapping, compile hooks, and UI wiring without spawning \texttt{motivoc}.
Compiler CTest and golden fixtures own language correctness end-to-end.
